\clearpage\thispagestyle{empty}
\pdfbookmark[0]{Sources for the review of Harding's star atlas}{harding1812-sources}
\begin{center}\textbf{Sources for the review of Harding's star atlas}\end{center}
\noindent This AI-translated review follows Carl Friedrich Gauss, \textit{Werke}, vol.~VI (Göttingen, 1874), pp.~541--542. It first appeared under Göttingen in \textit{Göttingische gelehrte Anzeigen}, no.~158, 3 October 1812, pp.~1569--1571.
\par\smallskip\noindent The collected text has no comma between Scorpius and Sagittarius. The first printing separates the names. The historical constellation names refer to the charts discussed in the review. Harding's later \href{https://www.davidrumsey.com/luna/servlet/detail/RUMSEY~8~1~372238~90138977:Title--Atlas-novus-coelestis-XXVII-}{1822 atlas} is available digitally through the David Rumsey Historical Map Collection.
\par\smallskip\noindent \href{https://gdz.sub.uni-goettingen.de/id/PPN236007467}{Collected text}. First printing, \href{https://images.p.sub.uni-goettingen.de/iiif/image/gdz:PPN319721507_1812:00001613/full/1650,/0/default.jpg}{p.~1569}, \href{https://images.p.sub.uni-goettingen.de/iiif/image/gdz:PPN319721507_1812:00001614/full/1650,/0/default.jpg}{p.~1570} and \href{https://images.p.sub.uni-goettingen.de/iiif/image/gdz:PPN319721507_1812:00001615/full/1650,/0/default.jpg}{p.~1571}.
\par\smallskip\noindent Translation and typesetting produced by AI.
\clearpage
